\documentclass[tikz,border=8pt]{standalone}
\usepackage{ctex}
\usepackage{amsmath}
\usepackage{amssymb}
\usepackage{pgfplots}
\pgfplotsset{compat=1.18}
\usetikzlibrary{calc,positioning,arrows.meta,patterns}

\begin{document}
\begin{tikzpicture}

%% ===== 信赖域方法示意图 =====
%% 等高线 + 当前点 x_k + 圆形信赖域 + 二次模型极小 vs 真实函数极小

%% --- 等高线：椭圆，模拟 f(x,y) = 2x^2 + 5*(y-0.5)^2 ---
%% 中心在 x^* = (0, 0.5)，但画时平移到 (0,0) 用于 standalone
%% 真实函数等高线，中心在 (0, 0.5)

%% 坐标系
\draw[->, gray!50, thin] (-3.5, 0) -- (3.5, 0) node[right, font=\scriptsize] {$x_1$};
\draw[->, gray!50, thin] (0, -2.5) -- (0, 3.2) node[above, font=\scriptsize] {$x_2$};

%% 等高线（椭圆，中心在真实极小 (0, 0.5)）
\foreach \r in {0.5, 1.0, 1.6, 2.2, 2.8} {
  \draw[gray!40, thin] (0, 0.5) ellipse ({\r*1.4} and {\r*0.9});
}
%% 真实极小值 x*
\fill[green!50!black] (0, 0.5) circle (3pt);
\node[right, font=\small, green!50!black] at (0.1, 0.5) {真实极小 $\mathbf{x}^*$};

%% 当前迭代点 x_k
\coordinate (Xk) at (-1.8, -0.8);
\fill[black] (Xk) circle (3pt);
\node[below left, font=\small] at (Xk) {$\mathbf{x}_k$};

%% 信赖域（圆形，半径 Delta_k）
\def\Delta{1.7}
\draw[blue!70!black, thick, dashed]
  (Xk) circle (\Delta)
  node[above right, xshift=0.8cm, yshift=0.8cm, font=\small, blue!70!black]
  {信赖域 $\|\mathbf{p}\| \leq \Delta_k$};

%% 标注半径
\draw[blue!60!black, <->, thin]
  (Xk) -- ++ ({-\Delta*0.707}, {\Delta*0.707})
  node[midway, above left, font=\scriptsize, blue!60!black] {$\Delta_k$};

%% 二次模型极小（在信赖域内，偏向真实极小方向）
%% 二次模型是 x_k 处的 Taylor 二阶展开，其无约束极小在 p^B = -H^{-1}g
%% 设 p^B 指向 (1.0, 0.9)（在域内），故二次模型极小在 x_k + p^B
\coordinate (ModelMin) at (-0.8, 0.1);  %% x_k + p^B（在信赖域内）
\fill[red!70!black] (ModelMin) circle (3pt);
\node[above right, font=\small, red!70!black] at (ModelMin)
  {二次模型极小 $\mathbf{x}_k + \mathbf{p}^B$};

%% 箭头：从 x_k 指向二次模型极小（牛顿步 p^B）
\draw[->, thick, red!70!black, >=Stealth]
  (Xk) -- (ModelMin)
  node[midway, right, font=\scriptsize, red!70!black] {$\mathbf{p}^B$};

%% Cauchy 点（沿梯度方向，在域边界上）
%% 梯度方向：从 x_k 指向梯度方向，设为 (-0.6, -1.0) 方向（归一化后）
%% Cauchy 点在域边界：x_k + Delta * g/|g|
\pgfmathsetmacro{\gx}{0.5}  %% 负梯度方向（下降）
\pgfmathsetmacro{\gy}{1.3}
\pgfmathsetmacro{\gn}{sqrt(\gx*\gx + \gy*\gy)}
\pgfmathsetmacro{\cpx}{-1.8 + \Delta*\gx/\gn}
\pgfmathsetmacro{\cpy}{-0.8 + \Delta*\gy/\gn}
\fill[orange!70!black] (\cpx, \cpy) circle (2.5pt);
\node[above, font=\small, orange!70!black] at (\cpx, \cpy)
  {Cauchy 点 $\mathbf{x}_k + \mathbf{p}^C$};

%% 箭头：从 x_k 指向 Cauchy 点
\draw[->, thick, orange!60!black, >=Stealth, dashed]
  (Xk) -- (\cpx, \cpy)
  node[midway, left, font=\scriptsize, orange!60!black] {$\mathbf{p}^C$};

%% 图例文字注释框
\node[draw=gray!40, fill=white, rounded corners=3pt,
      font=\scriptsize, inner sep=5pt, align=left, anchor=north east]
  at (3.3, 3.0) {
    \textbf{信赖域方法核心要素}\\[2pt]
    {\color{black}$\mathbf{x}_k$}：当前迭代点\\
    {\color{blue!70!black}圆形}：信赖域 $\|\mathbf{p}\|\leq\Delta_k$\\
    {\color{red!70!black}$\mathbf{x}_k+\mathbf{p}^B$}：二次模型极小\\
    {\color{orange!70!black}$\mathbf{x}_k+\mathbf{p}^C$}：Cauchy 点（沿梯度）\\
    {\color{green!50!black}$\mathbf{x}^*$}：真实函数极小
  };

%% 下降比标注
\node[draw=gray!30, fill=gray!5, rounded corners=3pt,
      font=\scriptsize, inner sep=4pt, align=left]
  at (2.0, -1.8) {
    下降比 $\rho_k = \dfrac{\text{ared}}{\text{pred}}$\\[2pt]
    $< \eta_1$：收缩 $\Delta_k$\\
    $\in [\eta_1,\eta_2)$：保持 $\Delta_k$\\
    $\geq \eta_2$：扩大 $\Delta_k$
  };

%% 标题
\node[font=\small\bfseries] at (0, -2.8)
  {信赖域方法：在可信球内同时确定步长和方向};

\end{tikzpicture}
\end{document}
